\documentclass[11pt]{article}
% ===========================================================================
%  Shared preamble for the AP Calculus AB exam-preparation worksheets.
%  Input right after \documentclass. Files beginning with "_" are not compiled
%  on their own. Navy and gold, to match the course's lesson pages.
% ===========================================================================
\usepackage[letterpaper,top=2.3cm,bottom=2.2cm,left=1.7cm,right=1.7cm,headheight=28pt,footskip=30pt]{geometry}
\usepackage{amsmath,amssymb}
\usepackage[T1]{fontenc}
\usepackage[utf8]{inputenc}
\usepackage{lmodern}
\usepackage{xcolor}
\usepackage[most]{tcolorbox}
\usepackage{enumitem}
\usepackage{fancyhdr}
\usepackage{array}
\usepackage{tikz}
\usepackage{pgfplots}
\pgfplotsset{compat=1.18}

\definecolor{navy}{HTML}{161E54}
\definecolor{ink}{HTML}{27357E}
\definecolor{gold}{HTML}{B98A00}
\definecolor{goldbg}{HTML}{FFFDF3}
\definecolor{inkbg}{HTML}{F6F7FD}
\definecolor{rule}{HTML}{C9CFF0}
\definecolor{okgreen}{HTML}{1B7A44}
\definecolor{okbg}{HTML}{F1F8F3}

\pagestyle{fancy}
\fancyhf{}
\fancyhead[L]{\footnotesize Name: \underline{\hspace{4cm}}}
\fancyhead[C]{\textbf{\color{navy}Integration Academy}\\[-2pt]{\scriptsize Calculus AB \textperiodcentered{} AP\textsuperscript{\textregistered} Exam Preparation}}
\fancyfoot[L]{\tiny\color{black!55}\parbox[t]{0.78\textwidth}{AP\textsuperscript{\textregistered} is a registered trademark of the College Board, which was not involved in the production of, and does not endorse, this material. \textcopyright{} Integration Academy.}}
\fancyfoot[R]{\scriptsize\thepage}
\renewcommand{\headrulewidth}{1.2pt}
\renewcommand{\headrule}{\hbox to\headwidth{\color{gold}\leaders\hrule height \headrulewidth\hfill}}
\renewcommand{\footrulewidth}{0pt}
\setlength{\parindent}{0pt}
\setlength{\parskip}{4pt}

% Title block: \wstitle{code}{title}{unit line}
\newcommand{\wstitle}[3]{\begin{center}{\small\bfseries\color{gold}WORKSHEET #1}\\[3pt]{\LARGE\bfseries\color{navy}#2}\\[3pt]{\small\color{black!60}#3}\end{center}\vspace{2pt}}
% Section heading
\newcommand{\sech}[1]{\par\vspace{10pt}{\large\bfseries\color{navy}#1}\par\nobreak\vspace{2pt}{\color{rule}\hrule height0.8pt}\vspace{7pt}}
% A level tag printed beside a group of questions
\newcommand{\level}[1]{\par\vspace{6pt}{\footnotesize\bfseries\color{gold}\MakeUppercase{#1}}\par\nobreak\vspace{2pt}}

\newtcolorbox{ideabox}{colframe=ink,colback=inkbg,coltitle=white,colbacktitle=ink,fonttitle=\bfseries\small,title={Key ideas},arc=2pt,boxrule=0.8pt,left=8pt,right=8pt}
\newtcolorbox{example}[1]{colframe=rule,colback=white,coltitle=navy,colbacktitle=inkbg,fonttitle=\bfseries\small,title={#1},arc=2pt,boxrule=0.8pt,borderline west={3pt}{0pt}{ink},left=10pt,right=8pt}
\newtcolorbox{frbox}[1]{colframe=gold,colback=goldbg,coltitle=white,colbacktitle=gold,fonttitle=\bfseries\small,title={#1},arc=2pt,boxrule=0.8pt,breakable,left=8pt,right=8pt}
\newtcolorbox{solbox}{colframe=okgreen,colback=okbg,coltitle=white,colbacktitle=okgreen,fonttitle=\bfseries\small,title={Worked solutions},arc=2pt,boxrule=0.8pt,breakable,left=8pt,right=8pt}
\newtcolorbox{scorebox}[1]{colframe=navy,colback=white,coltitle=white,colbacktitle=navy,fonttitle=\bfseries\small,title={#1},arc=2pt,boxrule=0.8pt,breakable,left=8pt,right=8pt}

\newcommand{\soln}{\par\smallskip\textbf{\color{ink}Solution.}\ }
\newcommand{\result}[1]{\par\smallskip\textbf{#1}}
% blank working space of a given height
\newcommand{\work}[1]{\par\vspace{#1}}
% a point value in the margin of a scoring guide
\newcommand{\pts}[1]{\hfill{\footnotesize\bfseries\color{gold}[#1]}}

% Coordinate grid for graph questions: \begin{apgraph}{xmin}{xmax}{ymin}{ymax} ... \end{apgraph}
\newenvironment{apgraph}[4]{\begin{center}\begin{tikzpicture}\begin{axis}[
  axis lines=middle, xmin=#1, xmax=#2, ymin=#3, ymax=#4,
  grid=both, grid style={black!12}, xtick distance=1, ytick distance=1,
  ticklabel style={font=\scriptsize}, xlabel={$x$}, ylabel={$y$},
  width=10cm, height=7.2cm, samples=80, clip=false,
  every axis plot/.append style={line width=1.1pt, color=ink}]}{\end{axis}\end{tikzpicture}\end{center}}
\newcommand{\closeddot}[2]{\addplot[only marks, mark=*, mark size=2.2pt, color=ink] coordinates {(#1,#2)};}
\newcommand{\opendot}[2]{\addplot[only marks, mark=*, mark size=2.2pt, color=ink, mark options={fill=white}] coordinates {(#1,#2)};}

% Multiple-choice options: four across, or stacked when the options are long.
\newcommand{\opts}[4]{\par\vspace{3pt}\begin{tabular}{@{}p{0.235\linewidth}p{0.235\linewidth}p{0.235\linewidth}p{0.235\linewidth}@{}}(A)\ #1 & (B)\ #2 & (C)\ #3 & (D)\ #4\end{tabular}\par\vspace{9pt}}
\newcommand{\optsv}[4]{\par\vspace{3pt}\begin{tabular}{@{}l@{\ \ }p{0.9\linewidth}@{}}(A) & #1\\(B) & #2\\(C) & #3\\(D) & #4\end{tabular}\par\vspace{9pt}}

\fancyhead[R]{\footnotesize\color{black!60}Worksheet 1.2}
\begin{document}

\wstitle{1.2}{Reading Limits off a Graph}{Unit 1: Limits and Continuity}

\begin{ideabox}
\begin{itemize}[leftmargin=*,itemsep=2pt,topsep=2pt]
  \item To read a limit, cover the point and follow the curve toward it from each side.
  \item A filled dot gives the \textbf{value} $f(c)$; it says nothing about the limit.
  \item A limit fails to exist at a jump, where the function is unbounded, or where it oscillates.
  \item For $\displaystyle\lim_{x\to c}f(g(x))$: find what $g(x)$ approaches \emph{and from which side}, then take that one-sided limit of $f$.
\end{itemize}
\end{ideabox}

\sech{Worked examples}
\begin{example}{Example 1 \textperiodcentered{} Limit and value at the same point}
The graph of $p$ is shown on the next page. Find $\displaystyle\lim_{x\to -1}p(x)$ and $p(-1)$.\soln From both sides the graph heads for the open circle at height 3. The filled dot at $x=-1$ is at height 1.\result{$\displaystyle\lim_{x\to -1}p(x)=3$ and $p(-1)=1$.}
\end{example}

\begin{example}{Example 2 \textperiodcentered{} A jump}
For the same function, find $\displaystyle\lim_{x\to 2^-}p(x)$, $\displaystyle\lim_{x\to 2^+}p(x)$ and $\displaystyle\lim_{x\to 2}p(x)$.\soln From the left the parabola comes down to the filled dot at height 0. From the right the line starts at the open circle at height 1.\result{$\displaystyle\lim_{x\to 2^-}p(x)=0$, $\displaystyle\lim_{x\to 2^+}p(x)=1$, and $\displaystyle\lim_{x\to 2}p(x)$ does not exist.}
\end{example}

\begin{example}{Example 3 \textperiodcentered{} A composition}
For the same function, find $\displaystyle\lim_{x\to 2^-}p\big(p(x)\big)$.\soln As $x\to2^-$, $p(x)=4-x^2\to0$ through values \emph{greater} than 0. So we need $\displaystyle\lim_{u\to 0^+}p(u)=4$.\result{$\displaystyle\lim_{x\to 2^-}p\big(p(x)\big)=4$.}
\end{example}

\newpage
\sech{Practice}
All questions on this page refer to the graph of $p$ below, defined on $[-4,5]$.
\begin{apgraph}{-5}{6}{-1.5}{5.5}
\addplot[domain=-4:-1]{x+4};
\addplot[domain=-1:2]{4-x^2};
\addplot[domain=2:5]{x-1};
\closeddot{-4}{0}\closeddot{-1}{1}\closeddot{2}{0}\closeddot{5}{4}
\opendot{-1}{3}\opendot{2}{1}
\node[font=\small, color=ink] at (axis cs:4.1,2.1) {$y=p(x)$};
\end{apgraph}
\level{Core}
\begin{enumerate}[leftmargin=*,itemsep=2pt]
  \item Find $\displaystyle\lim_{x\to 0}p(x)$.\work{1.2cm}
  \item Find $p(2)$.\work{1.2cm}
  \item Find $\displaystyle\lim_{x\to -1^+}p(x)$.\work{1.2cm}
  \item Find $\displaystyle\lim_{x\to 5^-}p(x)$.\work{1.2cm}
\end{enumerate}

\level{Exam level}
\begin{enumerate}[leftmargin=*,itemsep=2pt,start=5]
  \item Evaluate $\displaystyle\lim_{x\to -1}p(x)+p(2)$.\work{1.8cm}
  \item Evaluate $\displaystyle\lim_{x\to 1}\big[p(x)\big]^2$.\work{1.8cm}
  \item For which value(s) of $a$ in $(-4,5)$ does $\displaystyle\lim_{x\to a}p(x)$ exist but differ from $p(a)$?\work{1.8cm}
  \item At how many values of $x$ in $(-4,5)$ is $p$ discontinuous? List them.\work{1.8cm}
\end{enumerate}

\begin{frbox}{Question 9 \textperiodcentered{} Free-response practice}
Use the graph of $p$ on the previous page.
\begin{enumerate}[label=(\alph*),leftmargin=*,itemsep=2pt]
  \item Does $\displaystyle\lim_{x\to 2}p(x)$ exist? Justify your answer.\work{2.2cm}
  \item Find $\displaystyle\lim_{x\to 2^+}p\big(p(x)\big)$. Show your reasoning.\work{2.2cm}
  \item Does $\displaystyle\lim_{x\to 2}p\big(p(x)\big)$ exist? Justify, using part (b) and Example 3.\work{2.2cm}
\end{enumerate}
\end{frbox}

\newpage
\begin{solbox}
\begin{enumerate}[leftmargin=*,itemsep=4pt]
  \item $4$.
  \item $0$ (the filled dot).
  \item $3$.
  \item $4$.
  \item $3+0=3$.
  \item $p(x)\to3$ as $x\to1$, so the limit is $9$.
  \item $a=-1$ only: the limit is 3 but $p(-1)=1$. (At $x=2$ the limit does not exist.)
  \item Two: $x=-1$ (limit 3, value 1) and $x=2$ (a jump).
  \item \begin{enumerate}[label=(\alph*),leftmargin=*,itemsep=3pt]
    \item No. $\displaystyle\lim_{x\to 2^-}p(x)=0$ and $\displaystyle\lim_{x\to 2^+}p(x)=1$. Since the one-sided limits are not equal, the limit does not exist.
    \item As $x\to2^+$, $p(x)=x-1\to1$ through values greater than 1. Near $u=1$ the graph is the parabola, which is continuous there, so the limit is $p(1)=4-1=3$.
    \item No. $\displaystyle\lim_{x\to 2^-}p(p(x))=4$ (Example 3) and $\displaystyle\lim_{x\to 2^+}p(p(x))=3$. They are not equal, so the limit does not exist.
  \end{enumerate}
\end{enumerate}
\end{solbox}

\end{document}
